\documentclass[11pt,a4paper]{article}
\usepackage[utf8]{inputenc}
\usepackage[italian]{babel}
\usepackage{geometry}
\usepackage{graphicx}
\usepackage{hyperref}
\usepackage{listings}
\usepackage{xcolor}
\usepackage{amsmath}
\usepackage{enumitem}
\usepackage{booktabs}
\usepackage{longtable}
\usepackage{array}
\usepackage{tcolorbox}

\geometry{margin=2.5cm}

% Configurazione codice
\lstset{
    language=JavaScript,
    basicstyle=\ttfamily\small,
    keywordstyle=\color{blue}\bfseries,
    stringstyle=\color{red},
    commentstyle=\color{green!60!black},
    numberstyle=\tiny\color{gray},
    numbers=left,
    numbersep=5pt,
    frame=single,
    breaklines=true,
    showstringspaces=false,
    tabsize=2,
    escapeinside={(*@}{@*)}
}

% Box per evidenziare sezioni importanti
\newtcolorbox{importantbox}[1]{
    colback=red!5!white,
    colframe=red!75!black,
    title=#1,
    fonttitle=\bfseries
}

\newtcolorbox{warningbox}[1]{
    colback=orange!5!white,
    colframe=orange!75!black,
    title=#1,
    fonttitle=\bfseries
}

\newtcolorbox{tipbox}[1]{
    colback=blue!5!white,
    colframe=blue!75!black,
    title=#1,
    fonttitle=\bfseries
}

\newtcolorbox{successbox}[1]{
    colback=green!5!white,
    colframe=green!75!black,
    title=#1,
    fonttitle=\bfseries
}

\title{\textbf{Refactoring: Unificazione WorkLog e WorkEntry}\\
\large Documentazione Tecnica Completa del Refactoring}
\author{Documentazione Sistema}
\date{\today}

\begin{document}

\maketitle
\tableofcontents
\newpage

\section{Introduzione}

\subsection{Obiettivo del Refactoring}

Questo documento descrive in dettaglio il refactoring completo che ha unificato due entità separate (\texttt{WorkLog} e \texttt{WorkEntry}) in un'unica entità ibrida (\texttt{WorkLog}) per semplificare l'architettura e ridurre la duplicazione di codice.

\begin{importantbox}{Problema Risolto}
Prima del refactoring, il sistema aveva:
\begin{itemize}
    \item \textbf{WorkLog}: Gestiva solo log con orari (Timer) per il tracciamento delle ore
    \item \textbf{WorkEntry}: Gestiva solo note testuali (Journal) per il diario di lavoro
    \item \textbf{Duplicazione}: Due modelli, due servizi, due set di API per funzionalità simili
    \item \textbf{Complessità}: Logica di business duplicata e difficile da mantenere
\end{itemize}
\end{importantbox}

\subsection{Risultato Atteso}

Dopo il refactoring, il sistema ha:
\begin{itemize}
    \item \textbf{Un solo modello}: \texttt{WorkLog} che supporta sia Timer che Journal
    \item \textbf{Un solo servizio}: \texttt{WorkLogService} con logica ibrida
    \item \textbf{API unificate}: Un unico set di endpoint REST sotto \texttt{/api/worklogs}
    \item \textbf{UI semplificata}: Componenti React aggiornati per usare il modello unificato
\end{itemize}

\subsection{Architettura del Refactoring}

Il refactoring è stato eseguito in \textbf{4 fasi sequenziali}:

\begin{enumerate}
    \item \textbf{FASE 1: Database Refoundation (Model Layer)} - Modifica del modello Mongoose
    \item \textbf{FASE 2: Hybrid Business Logic (Service Layer)} - Aggiornamento della logica di business
    \item \textbf{FASE 3: API Exposure (Controller \& Routes)} - Consolidamento degli endpoint REST
    \item \textbf{FASE 4: UI Update (Frontend)} - Aggiornamento dei componenti React
\end{enumerate}

\section{FASE 1: Database Refoundation (Model Layer)}

\subsection{Obiettivo}

Unificare lo schema del database rendendo \texttt{WorkLog} in grado di gestire sia log con orari (Timer) che note senza orari (Journal).

\subsection{Modifiche al Modello WorkLog}

\subsubsection{Schema Prima del Refactoring}

Prima del refactoring, \texttt{WorkLog} aveva solo campi per il tracciamento temporale:

\begin{lstlisting}[caption=Schema WorkLog Originale]
const workLogSchema = new mongoose.Schema({
    user: { type: mongoose.Schema.Types.ObjectId, ref: 'User', required: true },
    projectId: { type: mongoose.Schema.Types.ObjectId, ref: 'Project', required: true },
    date: { type: String, required: true }, // YYYY-MM-DD
    startTime: { type: String, required: true }, // HH:mm
    endTime: { type: String, required: true }, // HH:mm
    durationMinutes: { type: Number, default: 0 },
    createdAt: { type: Date, default: Date.now },
    updatedAt: { type: Date, default: Date.now }
});
\end{lstlisting}

\subsubsection{Schema Dopo il Refactoring}

Dopo il refactoring, \texttt{WorkLog} supporta entrambi i casi d'uso:

\begin{lstlisting}[caption=Schema WorkLog Unificato]
const workLogSchema = new mongoose.Schema({
    user: { type: mongoose.Schema.Types.ObjectId, ref: 'User', required: true },
    projectId: { type: mongoose.Schema.Types.ObjectId, ref: 'Project', required: true },
    date: { type: String, required: true }, // YYYY-MM-DD
    
    // Campi per contenuto ricco (Journal)
    title: { type: String, required: true, maxlength: 200 },
    description: { type: String, maxlength: 5000 },
    tags: [{ type: String, trim: true }],
    mood: { 
        type: String, 
        enum: ['high', 'neutral', 'low'],
        default: undefined 
    },
    isBillable: { type: Boolean, default: true },
    
    // Campi temporali (Timer) - ORA OPZIONALI
    startTime: { type: String }, // HH:mm - OPZIONALE
    endTime: { type: String }, // HH:mm - OPZIONALE
    durationMinutes: { type: Number, default: 0 },
    
    createdAt: { type: Date, default: Date.now },
    updatedAt: { type: Date, default: Date.now }
});
\end{lstlisting}

\subsection{Modifiche Chiave}

\begin{enumerate}
    \item \textbf{Campi Aggiunti}:
    \begin{itemize}
        \item \texttt{title} (String, required): Titolo dell'attività
        \item \texttt{description} (String, max 5000): Note dettagliate
        \item \texttt{tags} ([String]): Tag per categorizzazione
        \item \texttt{mood} (String enum): Umore durante il lavoro
        \item \texttt{isBillable} (Boolean, default true): Indica se fatturabile
    \end{itemize}
    
    \item \textbf{Campi Resi Opzionali}:
    \begin{itemize}
        \item \texttt{startTime}: Rimosso \texttt{required: true}
        \item \texttt{endTime}: Rimosso \texttt{required: true}
    \end{itemize}
    
    \item \textbf{Index Aggiunti}:
    \begin{itemize}
        \item Text index su \texttt{title} e \texttt{description} per ricerca full-text
        \item Index su \texttt{tags} per filtri rapidi
    \end{itemize}
\end{enumerate}

\subsection{Pre-save Hook Aggiornato}

Il pre-save hook è stato modificato per gestire entrambi i casi:

\begin{lstlisting}[caption=Pre-save Hook Ibrido]
workLogSchema.pre('save', function() {
    // Se ci sono orari, calcola la durata
    if (this.startTime && this.endTime) {
        if (this.isModified('startTime') || this.isModified('endTime')) {
            const [startHour, startMin] = this.startTime.split(':').map(Number);
            const [endHour, endMin] = this.endTime.split(':').map(Number);
            
            let duration = (endHour * 60 + endMin) - (startHour * 60 + startMin);
            
            // Gestione lavoro oltre mezzanotte
            if (duration < 0) {
                duration += 24 * 60;
            }
            
            this.durationMinutes = duration;
        }
    } else {
        // Se non ci sono orari, durata = 0 (o valore manuale se presente)
        if (!this.isModified('durationMinutes')) {
            this.durationMinutes = 0;
        }
    }
    
    this.updatedAt = new Date();
});
\end{lstlisting}

\begin{tipbox}{Logica Ibrida}
Il hook distingue automaticamente tra:
\begin{itemize}
    \item \textbf{Timer}: Se \texttt{startTime} e \texttt{endTime} sono presenti, calcola la durata
    \item \textbf{Journal}: Se mancano gli orari, imposta \texttt{durationMinutes = 0}
\end{itemize}
\end{tipbox}

\subsection{Eliminazione WorkEntry}

Il modello \texttt{WorkEntry} è stato completamente eliminato:

\begin{lstlisting}[caption=Eliminazione File]
# File eliminato
server/models/WorkEntry.js
\end{lstlisting}

\begin{warningbox}{Breaking Change}
Tutti i dati esistenti in \texttt{WorkEntry} devono essere migrati a \texttt{WorkLog} prima di eliminare il modello. Questo documento assume che la migrazione sia già stata completata.
\end{warningbox}

\section{FASE 2: Hybrid Business Logic (Service Layer)}

\subsection{Obiettivo}

Aggiornare \texttt{WorkLogService} per gestire la logica ibrida che supporta sia Timer che Journal.

\subsection{Modifiche al WorkLogService}

\subsubsection{Configurazione BaseService}

Il servizio estende \texttt{BaseService} con configurazione aggiornata:

\begin{lstlisting}[caption=Configurazione WorkLogService]
class WorkLogService extends BaseService {
    constructor() {
        super(WorkLog, {
            searchFields: ['title', 'description', 'tags'],
            defaultSort: { date: -1, createdAt: -1 },
            populateFields: ['projectId'],
            // ... altre opzioni
        });
    }
}
\end{lstlisting}

\subsubsection{Metodo beforeCreate}

Il metodo \texttt{beforeCreate} è stato aggiornato per validare solo se necessario:

\begin{lstlisting}[caption=beforeCreate Ibrido]
async beforeCreate(tenantScope, data) {
    // 1. Validazione ownership progetto (sempre)
    if (data.projectId) {
        const owns = await tenantScope.owns(Project, data.projectId);
        if (!owns) {
            throw AppError.forbidden('Progetto non trovato o non accessibile');
        }
    }
    
    // 2. Validazione coerenza orari (solo se presenti)
    if (data.startTime && data.endTime) {
        const [startHour, startMin] = data.startTime.split(':').map(Number);
        const [endHour, endMin] = data.endTime.split(':').map(Number);
        
        if (isNaN(startHour) || isNaN(startMin) || 
            isNaN(endHour) || isNaN(endMin)) {
            throw AppError.validation('Formato orario non valido (usa HH:mm)');
        }
        
        // Il Model calcolerà la durata nel pre-save hook
    }
    
    // 3. Se non ci sono orari, è una nota/journal (OK)
    // durationMinutes sarà 0 o valore manuale se presente
    
    return data;
}
\end{lstlisting}

\subsubsection{Metodo beforeUpdate}

Il metodo \texttt{beforeUpdate} gestisce aggiornamenti parziali:

\begin{lstlisting}[caption=beforeUpdate Ibrido]
async beforeUpdate(tenantScope, documentId, data) {
    // 1. Validazione ownership progetto (se cambiato)
    if (data.projectId) {
        const owns = await tenantScope.owns(Project, data.projectId);
        if (!owns) {
            throw AppError.forbidden('Progetto non trovato o non accessibile');
        }
    }
    
    // 2. Validazione coerenza orari (se modificati)
    if (data.startTime !== undefined || data.endTime !== undefined) {
        // Prevenire aggiornamenti parziali di orari
        if ((data.startTime && !data.endTime) || 
            (!data.startTime && data.endTime)) {
            throw AppError.validation(
                'startTime e endTime devono essere entrambi presenti o assenti'
            );
        }
        
        // Se entrambi presenti, validare formato
        if (data.startTime && data.endTime) {
            const [startHour, startMin] = data.startTime.split(':').map(Number);
            const [endHour, endMin] = data.endTime.split(':').map(Number);
            
            if (isNaN(startHour) || isNaN(startMin) || 
                isNaN(endHour) || isNaN(endMin)) {
                throw AppError.validation('Formato orario non valido');
            }
        }
    }
    
    return data;
}
\end{lstlisting}

\subsubsection{Metodo create Override}

Il metodo \texttt{create} è stato sovrascritto per registrare attività appropriate:

\begin{lstlisting}[caption=create con Activity Tracking]
async create(tenantScope, data) {
    // Validazione e creazione tramite BaseService
    const created = await super.create(tenantScope, data);
    
    // Activity tracking basato sul tipo di entry
    if (data.startTime && data.endTime) {
        // Timer entry
        await activityService.recordActivity(tenantScope.tenantId, {
            type: 'WORK_SESSION_LOGGED',
            metadata: {
                workLogId: created._id,
                projectId: created.projectId,
                durationMinutes: created.durationMinutes,
            },
        });
    } else {
        // Journal entry
        await activityService.recordActivity(tenantScope.tenantId, {
            type: 'WORK_NOTE_CREATED',
            metadata: {
                workLogId: created._id,
                projectId: created.projectId,
            },
        });
    }
    
    return created;
}
\end{lstlisting}

\subsection{Nuovo Metodo: getWorkspaceFeed}

È stato aggiunto un nuovo metodo per la timeline del workspace:

\begin{lstlisting}[caption=getWorkspaceFeed]
async getWorkspaceFeed(tenantScope, filters = {}, options = {}) {
    const userId = this._getUserId(tenantScope);
    const queryFilter = { user: userId };
    
    // Filtro per progetto
    if (filters.projectId) {
        queryFilter.projectId = filters.projectId;
    }
    
    // Filtro per data singola
    if (filters.date) {
        queryFilter.date = filters.date;
    }
    
    // Filtro per range di date
    if (filters.startDate || filters.endDate) {
        queryFilter.date = {};
        if (filters.startDate) {
            queryFilter.date.$gte = filters.startDate;
        }
        if (filters.endDate) {
            queryFilter.date.$lte = filters.endDate;
        }
    }
    
    // Filtro per tags
    if (filters.tags) {
        const tagsArray = Array.isArray(filters.tags) 
            ? filters.tags 
            : [filters.tags];
        queryFilter.tags = { $in: tagsArray };
    }
    
    // Ordinamento timeline-friendly
    const sortOptions = { date: -1, createdAt: -1 };
    
    const query = this.Model.find(queryFilter)
        .sort(sortOptions)
        .populate(this.options.populateFields);
    
    if (options.limit) query.limit(options.limit);
    if (options.skip) query.skip(options.skip);
    
    return query.exec();
}
\end{lstlisting}

\subsection{Eliminazione WorkEntryService}

Il servizio \texttt{WorkEntryService} è stato eliminato:

\begin{lstlisting}[caption=Eliminazione Service]
# File eliminato
server/services/workEntryService.js
\end{lstlisting}

\subsection{Aggiornamento WorkspaceService}

Il \texttt{WorkspaceService} è stato aggiornato per rimuovere riferimenti a \texttt{WorkEntry}:

\begin{lstlisting}[caption=Aggiornamento WorkspaceService]
// PRIMA: Usava WorkEntry
const entriesCount = await WorkEntry.countDocuments({ 
    project: projectId 
});

// DOPO: Usa WorkLog
const entriesCount = await WorkLog.countDocuments({ 
    projectId: projectId 
});
\end{lstlisting}

\section{FASE 3: API Exposure (Controller \& Routes)}

\subsection{Obiettivo}

Consolidare tutti gli endpoint REST sotto \texttt{/api/worklogs} e rimuovere le vecchie route \texttt{/api/workspace/entries}.

\subsection{Modifiche al Controller}

\subsubsection{WorkLogController}

Il controller è stato aggiornato per accettare i nuovi campi:

\begin{lstlisting}[caption=createWorkLog Aggiornato]
exports.createWorkLog = asyncHandler(async (req, res) => {
    const {
        projectId,
        date,
        title,
        description,
        tags,
        mood,
        isBillable,
        startTime,
        endTime,
    } = req.body;
    
    // Validazione base
    if (!projectId || !date || !title) {
        throw AppError.validation('projectId, date e title sono obbligatori');
    }
    
    const data = {
        projectId,
        date,
        title,
        description,
        tags: Array.isArray(tags) ? tags : (tags ? [tags] : []),
        mood,
        isBillable: isBillable !== false, // default true
        startTime,
        endTime,
    };
    
    const workLog = await workLogService.create(req.tenantScope, data);
    
    res.status(201).json({
        success: true,
        data: workLog,
    });
});
\end{lstlisting}

\subsubsection{Nuovo Metodo: getWorkspaceFeed}

È stato aggiunto un nuovo metodo per la feed API:

\begin{lstlisting}[caption=getWorkspaceFeed Controller]
exports.getWorkspaceFeed = asyncHandler(async (req, res) => {
    const filters = {};
    
    if (req.query.projectId) {
        filters.projectId = req.query.projectId;
    }
    
    if (req.query.date) {
        filters.date = req.query.date;
    } else if (req.query.startDate || req.query.endDate) {
        filters.startDate = req.query.startDate;
        filters.endDate = req.query.endDate;
    }
    
    if (req.query.tags) {
        filters.tags = req.query.tags.includes(',') 
            ? req.query.tags.split(',').map(t => t.trim())
            : req.query.tags;
    }
    
    const options = {};
    if (req.query.limit) {
        options.limit = parseInt(req.query.limit);
    }
    if (req.query.skip) {
        options.skip = parseInt(req.query.skip);
    }
    
    const logs = await workLogService.getWorkspaceFeed(
        req.tenantScope, 
        filters, 
        options
    );
    
    res.json({
        success: true,
        data: logs,
    });
});
\end{lstlisting}

\subsection{Modifiche alle Routes}

\subsubsection{Routes Aggiornate}

Le route sono state consolidate:

\begin{lstlisting}[caption=Routes Unificate]
// POST /api/worklogs - Crea nuovo log (Timer o Journal)
router.post(
    '/worklogs',
    requireAuth,
    tenantContext(),
    workLogController.createWorkLog
);

// GET /api/worklogs/feed - Feed timeline con filtri
router.get(
    '/worklogs/feed',
    requireAuth,
    tenantContext(),
    workLogController.getWorkspaceFeed
);

// GET /api/worklogs/:id - Dettaglio singolo log
router.get(
    '/worklogs/:id',
    requireAuth,
    tenantContext(),
    workLogController.getWorkLog
);

// PUT /api/worklogs/:id - Aggiorna log
router.put(
    '/worklogs/:id',
    requireAuth,
    tenantContext(),
    workLogController.updateWorkLog
);

// DELETE /api/worklogs/:id - Elimina log
router.delete(
    '/worklogs/:id',
    requireAuth,
    tenantContext(),
    workLogController.deleteWorkLog
);
\end{lstlisting}

\subsubsection{Routes Rimosse}

Tutte le route vecchie sono state eliminate:

\begin{lstlisting}[caption=Routes Eliminate]
# Rimosse tutte le route:
# - POST /api/workspace/entries
# - GET /api/workspace/entries
# - GET /api/workspace/entries/:id
# - PUT /api/workspace/entries/:id
# - DELETE /api/workspace/entries/:id
# - GET /api/workspace/entries/timeline
# - GET /api/workspace/entries/by-month/:year/:month
# - GET /api/workspace/entries/by-project/:projectId
# - GET /api/workspace/entries/stats
\end{lstlisting}

\subsection{WorkspaceController Cleanup}

Il \texttt{WorkspaceController} è stato ripulito rimuovendo tutti i metodi relativi a \texttt{WorkEntry}:

\begin{lstlisting}[caption=Metodi Rimossi da WorkspaceController]
# Metodi eliminati:
# - getEntries
# - createEntry
# - updateEntry
# - deleteEntry
# - getTimeline
# - getEntriesByMonth
# - getEntriesByProject
# - getEntriesStats

# Metodi mantenuti (solo WorkProject e WorkTodo):
# - getProjects
# - createProject
# - updateProject
# - deleteProject
# - getTodos
# - createTodo
# - updateTodo
# - deleteTodo
\end{lstlisting}

\section{FASE 4: UI Update (Frontend)}

\subsection{Obiettivo}

Aggiornare tutti i componenti React per usare il modello unificato \texttt{WorkLog} invece di \texttt{WorkEntry}.

\subsection{Modifiche ai Tipi TypeScript}

\subsubsection{Nuovi Tipi}

Sono stati aggiunti nuovi tipi per \texttt{WorkLog}:

\begin{lstlisting}[caption=Tipi WorkLog]
export interface WorkLog {
    id: string;
    projectId: string;
    date: string; // YYYY-MM-DD
    title: string;
    description?: string;
    tags?: string[];
    mood?: 'high' | 'neutral' | 'low';
    isBillable?: boolean;
    startTime?: string; // HH:mm - OPZIONALE
    endTime?: string; // HH:mm - OPZIONALE
    durationMinutes?: number;
    createdAt: string;
    updatedAt: string;
}

export interface CreateWorkLogDTO {
    projectId: string;
    date: string;
    title: string;
    description?: string;
    tags?: string[];
    mood?: 'high' | 'neutral' | 'low';
    isBillable?: boolean;
    startTime?: string;
    endTime?: string;
}

export interface UpdateWorkLogDTO extends Partial<CreateWorkLogDTO> {}
\end{lstlisting}

\subsubsection{WorkEntry Deprecato}

Il tipo \texttt{WorkEntry} è stato marcato come deprecato:

\begin{lstlisting}[caption=WorkEntry Deprecato]
/**
 * @deprecated Usa WorkLog da '../tracker/type'
 */
export interface WorkEntry {
    // ... mantenuto per compatibilità temporanea
}
\end{lstlisting}

\subsection{Componenti Aggiornati}

\subsubsection{EntryCard.tsx}

Il componente è stato aggiornato per usare \texttt{WorkLog}:

\begin{lstlisting}[caption=EntryCard Aggiornato]
interface EntryCardProps {
    entry: WorkLog; // Cambiato da WorkEntry
    project: WorkProject | null;
    onClick?: () => void;
}

export const EntryCard = ({ entry, project, onClick }: EntryCardProps) => {
    const hasTimeTracking = !!(entry.startTime && entry.endTime);
    
    return (
        <motion.div>
            {/* Icona: Timer o Journal */}
            <div>
                {hasTimeTracking ? '⏱️' : '📝'}
            </div>
            
            {/* Mood invece di Category */}
            {entry.mood && (
                <span>
                    {MOOD_ICONS[entry.mood]} {MOOD_LABELS[entry.mood]}
                </span>
            )}
            
            {/* isBillable badge */}
            {entry.isBillable === false && (
                <span>Non fatturabile</span>
            )}
            
            {/* description invece di content */}
            {entry.description && (
                <p>{entry.description}</p>
            )}
            
            {/* Orari se presenti */}
            {hasTimeTracking && (
                <div>
                    {entry.startTime} - {entry.endTime}
                    {entry.durationMinutes > 0 && (
                        <span>{formatDuration(entry.durationMinutes)}</span>
                    )}
                </div>
            )}
        </motion.div>
    );
};
\end{lstlisting}

\subsubsection{EntryForm.tsx}

Il form è stato aggiornato per creare \texttt{WorkLog} (Journal):

\begin{lstlisting}[caption=EntryForm Aggiornato]
export const EntryForm = ({ onClose, entry, defaultProjectId }: EntryFormProps) => {
    const [title, setTitle] = useState(entry?.title || '');
    const [description, setDescription] = useState(entry?.description || '');
    const [tags, setTags] = useState<string[]>(entry?.tags || []);
    const [mood, setMood] = useState<'high' | 'neutral' | 'low' | undefined>(entry?.mood);
    const [isBillable, setIsBillable] = useState<boolean>(entry?.isBillable !== false);
    
    // Rimosso: category, content, duration
    
    const handleSubmit = async (e: React.FormEvent) => {
        const data: CreateWorkLogDTO = {
            projectId: finalProjectId,
            date: dateStr,
            title: title.trim(),
            description: description.trim() || undefined,
            tags: tags.length > 0 ? tags : undefined,
            mood: mood,
            isBillable: isBillable,
            // Non includere startTime/endTime per creare Journal
        };
        
        if (isEditing) {
            await updateEntry(entry.id, data);
        } else {
            await addEntry(data);
        }
    };
    
    return (
        <form>
            {/* Mood selector invece di Category */}
            <div>
                {MOOD_OPTIONS.map((m) => (
                    <button onClick={() => setMood(m.value)}>
                        {m.icon} {m.label}
                    </button>
                ))}
            </div>
            
            {/* isBillable checkbox */}
            <div>
                <input
                    type="checkbox"
                    checked={isBillable}
                    onChange={(e) => setIsBillable(e.target.checked)}
                />
                Lavoro fatturabile
            </div>
        </form>
    );
};
\end{lstlisting}

\subsubsection{EntryTimeline.tsx}

La timeline è stata aggiornata per usare \texttt{WorkLog[]}:

\begin{lstlisting}[caption=EntryTimeline Aggiornato]
interface EntryTimelineProps {
    entries: WorkLog[]; // Cambiato da WorkEntry[]
    loading: boolean;
    projects: WorkProject[];
    onEntryClick?: (entry: WorkLog) => void;
}

export const EntryTimeline = ({ entries, projects, onEntryClick }: EntryTimelineProps) => {
    const groupedEntries = useMemo(() => {
        const groups: Record<string, WorkLog[]> = {};
        
        entries.forEach((entry) => {
            if (!groups[entry.date]) {
                groups[entry.date] = [];
            }
            groups[entry.date].push(entry);
        });
        
        return Object.entries(groups)
            .sort(([a], [b]) => b.localeCompare(a))
            .map(([date, entries]) => ({ 
                date, 
                entries: entries.sort((a, b) => 
                    new Date(b.createdAt).getTime() - new Date(a.createdAt).getTime()
                )
            }));
    }, [entries]);
    
    const getProject = (entry: WorkLog): WorkProject | null => {
        return projects.find((p) => {
            const pid = p.id || (p as any)._id;
            return String(pid) === String(entry.projectId);
        }) || null;
    };
    
    // ... rendering
};
\end{lstlisting}

\subsubsection{EntryDetailModal.tsx}

Il modal è stato aggiornato per mostrare i nuovi campi:

\begin{lstlisting}[caption=EntryDetailModal Aggiornato]
export const EntryDetailModal = ({
    isOpen,
    entry,
    project,
    onClose,
    onEdit,
    onDelete,
}: EntryDetailModalProps) => {
    if (!entry) return null;
    
    const hasTimeTracking = !!(entry.startTime && entry.endTime);
    
    return (
        <motion.div>
            {/* Header con tipo */}
            <div>
                {hasTimeTracking ? '⏱️ Timer' : '📝 Journal'}
                {entry.mood && (
                    <span>
                        {MOOD_ICONS[entry.mood]} {MOOD_LABELS[entry.mood]}
                    </span>
                )}
                {entry.isBillable === false && (
                    <span>Non fatturabile</span>
                )}
            </div>
            
            {/* Contenuto */}
            <div>
                {entry.description && (
                    <div>
                        <h3>Note</h3>
                        <p>{entry.description}</p>
                    </div>
                )}
                
                {/* Orari se presenti */}
                {hasTimeTracking && (
                    <div>
                        {entry.startTime} - {entry.endTime}
                        {entry.durationMinutes > 0 && (
                            <span>{formatDuration(entry.durationMinutes)}</span>
                        )}
                    </div>
                )}
            </div>
        </motion.div>
    );
};
\end{lstlisting}

\subsection{Context Aggiornato}

\subsubsection{WorkspaceContext.tsx}

Il context è stato aggiornato per usare \texttt{WorkLog}:

\begin{lstlisting}[caption=WorkspaceContext Aggiornato]
interface WorkspaceContextType {
    // Entries (ora WorkLog)
    entries: WorkLog[]; // Cambiato da WorkEntry[]
    entriesLoading: boolean;
    entriesError: string | null;
    refreshEntries: () => Promise<void>;
    addEntry: (data: CreateWorkLogDTO) => Promise<void>; // Cambiato da CreateWorkEntryDTO
    updateEntry: (id: string, data: UpdateWorkLogDTO) => Promise<void>; // Cambiato da UpdateWorkEntryDTO
    deleteEntry: (id: string) => Promise<void>;
}

export const WorkspaceProvider = ({ children }: { children: ReactNode }) => {
    const [entries, setEntries] = useState<WorkLog[]>([]); // Cambiato da WorkEntry[]
    
    const refreshEntries = useCallback(async () => {
        // Usa il nuovo endpoint feed
        const data = await workspaceEntriesService.getAll();
        setEntries(data as WorkLog[]);
    }, [isAuthenticated]);
    
    const addEntry = useCallback(async (data: CreateWorkLogDTO) => {
        const created = await workspaceEntriesService.create(data as any);
        setEntries((prev) => [...prev, created as WorkLog]);
    }, []);
    
    // ... altri metodi
};
\end{lstlisting}

\subsection{Service Frontend Aggiornato}

\subsubsection{workspaceService.ts}

Il service frontend è stato aggiornato per normalizzare i dati:

\begin{lstlisting}[caption=workspaceService Aggiornato]
export const workspaceEntriesService = {
    async getAll(filters?: {
        projectId?: string;
        date?: string;
        startDate?: string;
        endDate?: string;
        tags?: string | string[];
        limit?: number;
    }): Promise<any[]> {
        const response = await apiClient.get<any[]>(
            `/worklogs/feed${queryString ? `?${queryString}` : ''}`
        );
        const data = response.data || [];
        
        // Normalizza i dati da WorkLog
        return data.map((item: any) => ({
            id: String(item.id || item._id),
            projectId: item.projectId || item.project,
            date: item.date,
            title: item.title,
            description: item.description,
            tags: Array.isArray(item.tags) ? item.tags : [],
            mood: item.mood,
            isBillable: item.isBillable !== false,
            startTime: item.startTime,
            endTime: item.endTime,
            durationMinutes: item.durationMinutes || 0,
            createdAt: item.createdAt || new Date().toISOString(),
            updatedAt: item.updatedAt || new Date().toISOString(),
        }));
    },
    
    async create(data: CreateWorkLogDTO | any): Promise<any> {
        const workLogData: any = {
            projectId: data.projectId || data.project,
            date: data.date,
            title: data.title,
            description: data.description || data.content,
            tags: data.tags || [],
            mood: data.mood,
            isBillable: data.isBillable !== false,
        };
        
        const response = await apiClient.post<any>('/worklogs', workLogData);
        // ... normalizzazione
    },
    
    // ... altri metodi
};
\end{lstlisting}

\subsection{Pagine Aggiornate}

\subsubsection{WorkspacePage.tsx}

La pagina principale è stata aggiornata:

\begin{lstlisting}[caption=WorkspacePage Aggiornato]
export const WorkspacePage = () => {
    const [selectedEntry, setSelectedEntry] = useState<WorkLog | null>(null); // Cambiato da WorkEntry
    const [editingEntry, setEditingEntry] = useState<WorkLog | undefined>(); // Cambiato da WorkEntry
    
    // Rimosso: selectedCategory (non più necessario)
    
    const filteredEntries = useMemo(() => {
        let filtered = [...entries];
        
        if (selectedProject) {
            filtered = filtered.filter((e) => {
                return String(e.projectId) === String(selectedProject);
            });
        }
        
        // Ricerca su title, description, tags
        if (searchQuery.trim()) {
            const query = searchQuery.toLowerCase();
            filtered = filtered.filter((e) => {
                const titleMatch = e.title.toLowerCase().includes(query);
                const descriptionMatch = e.description?.toLowerCase().includes(query);
                const tagMatch = e.tags?.some((tag) => tag.toLowerCase().includes(query));
                return titleMatch || descriptionMatch || tagMatch;
            });
        }
        
        return filtered.sort((a, b) => {
            const dateCompare = b.date.localeCompare(a.date);
            if (dateCompare !== 0) return dateCompare;
            return new Date(b.createdAt).getTime() - new Date(a.createdAt).getTime();
        });
    }, [entries, selectedProject, searchQuery]);
    
    // Rimosso: Filtro per categoria
};
\end{lstlisting}

\subsubsection{ProjectDetailPage.tsx}

La pagina dettaglio progetto è stata aggiornata:

\begin{lstlisting}[caption=ProjectDetailPage Aggiornato]
export const ProjectDetailPage = () => {
    const [selectedEntry, setSelectedEntry] = useState<WorkLog | null>(null); // Cambiato da WorkEntry
    
    // Rimosso: selectedCategory
    
    const projectEntries = useMemo(() => {
        if (!project) return [];
        
        const projectId = project.id || (project as any)._id;
        let filtered = entries.filter((e) => {
            return String(e.projectId) === String(projectId);
        });
        
        // Ricerca su title, description, tags
        if (searchQuery.trim()) {
            // ... logica ricerca
        }
        
        return filtered.sort((a, b) => {
            const dateCompare = b.date.localeCompare(a.date);
            if (dateCompare !== 0) return dateCompare;
            return new Date(b.createdAt).getTime() - new Date(a.createdAt).getTime();
        });
    }, [project, entries, searchQuery]);
    
    const stats = {
        totalEntries: projectEntries.length,
        totalDuration: projectEntries.reduce((sum, e) => sum + (e.durationMinutes || 0), 0),
        withTimeTracking: projectEntries.filter((e) => e.startTime && e.endTime).length,
        journalEntries: projectEntries.filter((e) => !e.startTime || !e.endTime).length,
    };
    
    // Rimosso: Filtro per categoria
};
\end{lstlisting}

\section{Breaking Changes e Migrazioni}

\subsection{Breaking Changes}

\begin{importantbox}{Breaking Changes}
I seguenti cambiamenti rompono la compatibilità con il codice esistente:
\begin{enumerate}
    \item \textbf{Modello WorkEntry eliminato}: Tutti i riferimenti a \texttt{WorkEntry} devono essere aggiornati
    \item \textbf{API Routes cambiate}: \texttt{/api/workspace/entries/*} non esiste più
    \item \textbf{Campi rinominati}: \texttt{content} → \texttt{description}, \texttt{project} → \texttt{projectId}
    \item \textbf{Campi rimossi}: \texttt{category}, \texttt{templateData}, \texttt{duration} (ora \texttt{durationMinutes})
    \item \textbf{Campi aggiunti}: \texttt{title} (required), \texttt{mood}, \texttt{isBillable}
\end{enumerate}
\end{importantbox}

\subsection{Migrazione Dati}

Se ci sono dati esistenti in \texttt{WorkEntry}, devono essere migrati:

\begin{lstlisting}[caption=Script Migrazione Dati]
// Esempio script migrazione (da eseguire manualmente)
async function migrateWorkEntriesToWorkLogs() {
    const entries = await WorkEntry.find({});
    
    for (const entry of entries) {
        const workLog = new WorkLog({
            user: entry.user,
            projectId: entry.project, // Converti riferimento
            date: entry.date,
            title: entry.title,
            description: entry.content, // content → description
            tags: entry.tags || [],
            mood: undefined, // Non presente in WorkEntry
            isBillable: true, // Default
            // startTime/endTime non presenti in WorkEntry
            durationMinutes: entry.duration || 0,
            createdAt: entry.createdAt,
            updatedAt: entry.updatedAt,
        });
        
        await workLog.save();
    }
    
    // Dopo migrazione, eliminare WorkEntry
    await WorkEntry.deleteMany({});
}
\end{lstlisting}

\subsection{Compatibilità Temporanea}

Per facilitare la migrazione, alcuni servizi supportano ancora formati legacy:

\begin{lstlisting}[caption=Supporto Legacy]
// workspaceService.ts supporta ancora:
const workLogData = {
    projectId: data.projectId || data.project, // Supporta entrambi
    description: data.description || data.content, // Supporta entrambi
    // ...
};
\end{lstlisting}

\begin{warningbox}{Deprecazione}
Il supporto legacy verrà rimosso in una versione futura. Aggiornare il codice il prima possibile.
\end{warningbox}

\section{Testing e Validazione}

\subsection{Test da Eseguire}

\begin{enumerate}
    \item \textbf{Creazione Timer Entry}:
    \begin{itemize}
        \item Verificare che con \texttt{startTime} e \texttt{endTime} la durata venga calcolata correttamente
        \item Verificare gestione lavoro oltre mezzanotte
        \item Verificare activity tracking \texttt{WORK_SESSION_LOGGED}
    \end{itemize}
    
    \item \textbf{Creazione Journal Entry}:
    \begin{itemize}
        \item Verificare che senza orari, \texttt{durationMinutes = 0}
        \item Verificare che tutti i campi opzionali funzionino
        \item Verificare activity tracking \texttt{WORK_NOTE_CREATED}
    \end{itemize}
    
    \item \textbf{Feed API}:
    \begin{itemize}
        \item Verificare filtri per \texttt{projectId}, \texttt{date}, \texttt{tags}
        \item Verificare ordinamento per data e \texttt{createdAt}
        \item Verificare paginazione con \texttt{limit} e \texttt{skip}
    \end{itemize}
    
    \item \textbf{UI Components}:
    \begin{itemize}
        \item Verificare che \texttt{EntryCard} mostri correttamente Timer vs Journal
        \item Verificare che \texttt{EntryForm} crei correttamente Journal entries
        \item Verificare che \texttt{EntryTimeline} raggruppi correttamente per data
    \end{itemize}
\end{enumerate}

\subsection{Validazione Backend}

\begin{lstlisting}[caption=Test Backend]
// Test creazione Timer
const timerEntry = await workLogService.create(tenantScope, {
    projectId: '...',
    date: '2024-01-15',
    title: 'Sviluppo feature',
    startTime: '09:00',
    endTime: '17:00',
});
// Verifica: durationMinutes === 480

// Test creazione Journal
const journalEntry = await workLogService.create(tenantScope, {
    projectId: '...',
    date: '2024-01-15',
    title: 'Note riunione',
    description: 'Discussione su architettura',
    tags: ['meeting'],
    mood: 'high',
    isBillable: false,
});
// Verifica: durationMinutes === 0, startTime/endTime undefined
\end{lstlisting}

\subsection{Validazione Frontend}

\begin{lstlisting}[caption=Test Frontend]
// Test EntryForm
const formData: CreateWorkLogDTO = {
    projectId: 'project123',
    date: '2024-01-15',
    title: 'Nuova entry',
    description: 'Descrizione',
    tags: ['tag1', 'tag2'],
    mood: 'high',
    isBillable: true,
};

await addEntry(formData);
// Verifica: Entry creata senza orari (Journal)
\end{lstlisting}

\section{Riepilogo e Conclusioni}

\subsection{Riepilogo Modifiche}

\begin{table}[h]
\centering
\begin{tabular}{lcc}
\toprule
\textbf{Componente} & \textbf{File Modificati} & \textbf{File Eliminati} \\
\midrule
Model Layer & 1 (WorkLog.js) & 1 (WorkEntry.js) \\
Service Layer & 2 (workLogService.js, workspaceService.js) & 1 (workEntryService.js) \\
Controller Layer & 1 (workLogController.js) & - \\
Routes & 1 (api.js) & - \\
WorkspaceController & 1 (workspaceController.js) & - \\
Frontend Components & 6 & - \\
Frontend Services & 1 (workspaceService.ts) & - \\
Frontend Context & 1 (WorkspaceContext.tsx) & - \\
Frontend Pages & 2 (WorkspacePage.tsx, ProjectDetailPage.tsx) & - \\
\bottomrule
\end{tabular}
\caption{Riepilogo Modifiche}
\end{table}

\subsection{Benefici Ottenuti}

\begin{successbox}{Benefici}
\begin{itemize}
    \item \textbf{Riduzione Complessità}: Un solo modello invece di due
    \item \textbf{Manutenibilità}: Codice più semplice da mantenere
    \item \textbf{Flessibilità}: Supporto nativo per Timer e Journal nello stesso modello
    \item \textbf{API Unificate}: Un solo set di endpoint REST
    \item \textbf{UI Semplificata}: Componenti più semplici e coerenti
\end{itemize}
\end{successbox}

\subsection{Prossimi Passi}

\begin{enumerate}
    \item \textbf{Migrazione Dati}: Eseguire migrazione di eventuali dati esistenti da \texttt{WorkEntry} a \texttt{WorkLog}
    \item \textbf{Rimozione Legacy}: Rimuovere supporto legacy nei servizi dopo migrazione completa
    \item \textbf{Testing Completo}: Eseguire test end-to-end su tutte le funzionalità
    \item \textbf{Documentazione Utente}: Aggiornare documentazione utente se necessario
\end{enumerate}

\section{Appendice: Codice Completo}

\subsection{WorkLog Model Completo}

\begin{lstlisting}[caption=server/models/WorkLog.js - Versione Finale]
const mongoose = require('mongoose');
const multiTenancyPlugin = require('../plugins/multiTenancy');

const workLogSchema = new mongoose.Schema({
    user: {
        type: mongoose.Schema.Types.ObjectId,
        ref: 'User',
        required: true,
        index: true,
    },
    projectId: {
        type: mongoose.Schema.Types.ObjectId,
        ref: 'Project',
        required: true,
        index: true,
    },
    date: {
        type: String,
        required: true,
        match: /^\d{4}-\d{2}-\d{2}$/,
        index: true,
    },
    
    // Campi per contenuto ricco (Journal)
    title: {
        type: String,
        required: true,
        maxlength: 200,
        trim: true,
    },
    description: {
        type: String,
        maxlength: 5000,
        trim: true,
    },
    tags: [{
        type: String,
        trim: true,
    }],
    mood: {
        type: String,
        enum: ['high', 'neutral', 'low'],
        default: undefined,
    },
    isBillable: {
        type: Boolean,
        default: true,
    },
    
    // Campi temporali (Timer) - OPZIONALI
    startTime: {
        type: String,
        match: /^([0-1][0-9]|2[0-3]):[0-5][0-9]$/,
    },
    endTime: {
        type: String,
        match: /^([0-1][0-9]|2[0-3]):[0-5][0-9]$/,
    },
    durationMinutes: {
        type: Number,
        default: 0,
        min: 0,
    },
    
    createdAt: {
        type: Date,
        default: Date.now,
    },
    updatedAt: {
        type: Date,
        default: Date.now,
    },
}, {
    timestamps: false, // Gestiamo manualmente
});

// Index per ricerca
workLogSchema.index({ title: 'text', description: 'text' });
workLogSchema.index({ tags: 1 });
workLogSchema.index({ date: -1, createdAt: -1 });

// Pre-save hook
workLogSchema.pre('save', function() {
    if (this.startTime && this.endTime) {
        if (this.isModified('startTime') || this.isModified('endTime')) {
            const [startHour, startMin] = this.startTime.split(':').map(Number);
            const [endHour, endMin] = this.endTime.split(':').map(Number);
            
            let duration = (endHour * 60 + endMin) - (startHour * 60 + startMin);
            
            if (duration < 0) {
                duration += 24 * 60;
            }
            
            this.durationMinutes = duration;
        }
    } else {
        if (!this.isModified('durationMinutes')) {
            this.durationMinutes = 0;
        }
    }
    
    this.updatedAt = new Date();
});

// Plugin multi-tenancy
workLogSchema.plugin(multiTenancyPlugin);

module.exports = mongoose.model('WorkLog', workLogSchema);
\end{lstlisting}

\subsection{WorkLogService Metodi Chiave}

\begin{lstlisting}[caption=server/services/workLogService.js - Metodi Principali]
class WorkLogService extends BaseService {
    constructor() {
        super(WorkLog, {
            searchFields: ['title', 'description', 'tags'],
            defaultSort: { date: -1, createdAt: -1 },
            populateFields: ['projectId'],
        });
    }
    
    async beforeCreate(tenantScope, data) {
        // Validazione ownership progetto
        if (data.projectId) {
            const owns = await tenantScope.owns(Project, data.projectId);
            if (!owns) {
                throw AppError.forbidden('Progetto non trovato o non accessibile');
            }
        }
        
        // Validazione orari (solo se presenti)
        if (data.startTime && data.endTime) {
            // Validazione formato...
        }
        
        return data;
    }
    
    async create(tenantScope, data) {
        const created = await super.create(tenantScope, data);
        
        // Activity tracking
        if (data.startTime && data.endTime) {
            await activityService.recordActivity(tenantScope.tenantId, {
                type: 'WORK_SESSION_LOGGED',
                metadata: { workLogId: created._id, ... },
            });
        } else {
            await activityService.recordActivity(tenantScope.tenantId, {
                type: 'WORK_NOTE_CREATED',
                metadata: { workLogId: created._id, ... },
            });
        }
        
        return created;
    }
    
    async getWorkspaceFeed(tenantScope, filters = {}, options = {}) {
        // ... implementazione completa
    }
}
\end{lstlisting}

\section{Storico Modifiche}

\begin{table}[h]
\centering
\begin{tabular}{llp{8cm}}
\toprule
\textbf{Data} & \textbf{Autore} & \textbf{Modifica} \\
\midrule
\today & Sistema & Documentazione completa refactoring unificazione WorkLog/WorkEntry \\
\bottomrule
\end{tabular}
\caption{Storico Modifiche Documento}
\end{table}

\end{document}
